% IDIOMA E TEXTO GERAL ------------------

\usepackage[brazil]{babel}
\usepackage[none]{hyphenat}

% LAYOUT GERAL --------------------------

\usepackage{fancyhdr}
\usepackage{multicol}
\usepackage{parskip}
\usepackage[
a4paper,
top=2cm,
bottom=1.4cm,
left=1.1cm,
right=1.1cm
]{geometry}

% CORES E VISUAL ------------------------

\usepackage{tcolorbox}
\usepackage{tikz}

% IMAGENS E LINKS -----------------------

\usepackage{graphicx}
\usepackage{hyperref}

% FONTE E TIPOGRAFIA --------------------

\usepackage{fontspec}

\setmainfont{Poppins}[
  Path = fonts/poppins/,
  Extension = .ttf,
  UprightFont = Poppins-Regular,
  BoldFont = Poppins-Bold,
  ItalicFont = Poppins-Italic,
  BoldItalicFont = Poppins-BoldItalic
]

\newfontfamily\opensansfont{Open Sans}[
  Path = fonts/openSans/,
  Extension = .ttf,
  UprightFont = OpenSans-Regular,
  BoldFont = OpenSans-Bold,
  ItalicFont = OpenSans-Italic,
  BoldItalicFont = OpenSans-BoldItalic
]
 
\usepackage{titlesec}
\usepackage{enumitem}